\documentclass{standalone}
\usepackage{circuitikz}
\usetikzlibrary{calc}
\definecolor{circuitnotemark}{HTML}{FE00FE}
\begin{document}
\begin{circuitikz}[american, line width=0.8pt]
\ctikzset{bipoles/length=1.2cm}
\foreach \x in {0,2,4,6,8,10,12,14,16,18,20} {\foreach \y in {0,-2} {\fill[gray, opacity=0.35] (\x,\y) circle (1.1pt);}}
\node[gray, font=\scriptsize] at (0,0.84) {1};
\node[gray, font=\scriptsize] at (2,0.84) {2};
\node[gray, font=\scriptsize] at (4,0.84) {3};
\node[gray, font=\scriptsize] at (6,0.84) {4};
\node[gray, font=\scriptsize] at (8,0.84) {5};
\node[gray, font=\scriptsize] at (10,0.84) {6};
\node[gray, font=\scriptsize] at (12,0.84) {7};
\node[gray, font=\scriptsize] at (14,0.84) {8};
\node[gray, font=\scriptsize] at (16,0.84) {9};
\node[gray, font=\scriptsize] at (18,0.84) {10};
\node[gray, font=\scriptsize] at (20,0.84) {11};
\node[gray, font=\scriptsize, anchor=east] at (-0.84,0) {a};
\node[gray, font=\scriptsize, anchor=east] at (-0.84,-2) {b};
\coordinate (b2) at (2,-2);
\coordinate (b5) at (8,-2);
\coordinate (b8) at (14,-2);
\coordinate (b11) at (20,-2);
\node[npn] (part-Q1) at (b2) {}; % line 3
\node[npn, rotate=-90] (part-Q2) at (b5) {}; % line 4
\node[npn, rotate=-180] (part-Q3) at (b8) {}; % line 5
\node[npn, rotate=-270] (part-Q4) at (b11) {}; % line 6
\path (1.238,-6.593) rectangle (2.762,-5.407); % line 8
\node[anchor=west, inner sep=0, circuitnotemark, font=\large] at (2,-6) {X}; % line 8
\path (6.73,-6.593) rectangle (9.27,-5.407); % line 9
\node[anchor=west, inner sep=0, circuitnotemark, font=\large] at (8,-6) {X}; % line 9
\path (12.603,-6.593) rectangle (15.397,-5.407); % line 10
\node[anchor=west, inner sep=0, circuitnotemark, font=\large] at (14,-6) {X}; % line 10
\path (18.476,-6.593) rectangle (21.524,-5.407); % line 11
\node[anchor=west, inner sep=0, circuitnotemark, font=\large] at (20,-6) {X}; % line 11
\path (24,-7.895) rectangle (34.499,0.593); % line 12
\node[anchor=west, inner sep=0, circuitnotemark, font=\large] at (24,0) {X}; % line 12
\node[anchor=west, inner sep=0, circuitnotemark, font=\large] at (24,-0.487) {X}; % line 12
\node[anchor=west, inner sep=0, circuitnotemark, font=\large] at (24,-0.974) {X}; % line 12
\node[anchor=west, inner sep=0, circuitnotemark, font=\large] at (24,-1.46) {X}; % line 12
\node[anchor=west, inner sep=0, circuitnotemark, font=\large] at (24,-1.947) {X}; % line 12
\node[anchor=west, inner sep=0, circuitnotemark, font=\large] at (24,-2.434) {X}; % line 12
\node[anchor=west, inner sep=0, circuitnotemark, font=\large] at (24,-2.921) {X}; % line 12
\node[anchor=west, inner sep=0, circuitnotemark, font=\large] at (24,-3.408) {X}; % line 12
\node[anchor=west, inner sep=0, circuitnotemark, font=\large] at (24,-3.895) {X}; % line 12
\node[anchor=west, inner sep=0, circuitnotemark, font=\large] at (24,-4.381) {X}; % line 12
\node[anchor=west, inner sep=0, circuitnotemark, font=\large] at (24,-4.868) {X}; % line 12
\node[anchor=west, inner sep=0, circuitnotemark, font=\large] at (24,-5.355) {X}; % line 12
\node[anchor=west, inner sep=0, circuitnotemark, font=\large] at (24,-5.842) {X}; % line 12
\node[anchor=west, inner sep=0, circuitnotemark, font=\large] at (24,-6.329) {X}; % line 12
\node[anchor=west, inner sep=0, circuitnotemark, font=\large] at (24,-6.816) {X}; % line 12
\node[anchor=west, inner sep=0, circuitnotemark, font=\large] at (24,-7.302) {X}; % line 12
\node[anchor=south west, inner sep=2pt, circuitnotemark, font=\LARGE\bfseries] (circuittitle) at (current bounding box.north west) {X};
\path (circuittitle.south west) rectangle ++(5.578,0.853);
\end{circuitikz}
\end{document}
